\section*{第 \ref{ch:b1:sinusoidal-impedance} 章——正弦稳态与阻抗}

\begin{solution}{exo:b1:sinusoidal-impedance:1}
$\underline U = 5\eu^{j\pi/3} = 2.5 + 4.33j$ (V)。$3\cos\omega t +
4\sin\omega t$：$3 + 4\eu^{-j\pi/2} = 3 - 4j$，模为 $5$，辐角为
$-\arctan(4/3) = \qty{-0.927}{rad}$，故为 $5\cos(\omega t - 0.927)$。
\end{solution}

\begin{solution}{exo:b1:sinusoidal-impedance:2}
$R$：两处均为 \qty{100}{\ohm}。$L\omega$：\qty{31.4}{\ohm}（\qty{50}{Hz}），
\qty{628}{\ohm}（\qty{1}{kHz}）。$1/C\omega$：\qty{318}{\ohm}，\qty{15.9}{\ohm}。
\qty{50}{Hz} 时的串联 $RL$：$\abs{\underline Z} = \sqrt{100^2 + 31.4^2} =
\qty{105}{\ohm}$，相位 $\arctan(0.314) = 17^\circ$。
\end{solution}

\begin{solution}{exo:b1:sinusoidal-impedance:3}
$U = 230\sqrt2 = \qty{325}{V}$。$I_{\mathrm{rms}} = 2000/230 = \qty{8.7}{A}$，
峰值为 \qty{12.3}{A}。由方均根值的定义，$\langle Ri^2\rangle = R\langle i^2\rangle =
RI_{\mathrm{rms}}^2$。
\end{solution}

\begin{solution}{exo:b1:sinusoidal-impedance:4}
$f_0 = 1/(2\pi\sqrt{LC}) = \qty{5.03}{kHz}$；$Q = \sqrt{L/C}/R = 3.16$；
$\Delta f = f_0/Q = \qty{1.6}{kHz}$；$I_{\max} = E/R = \qty{10}{mA}$。
\end{solution}

\begin{solution}{exo:b1:sinusoidal-impedance:5}
比值为 $1/\sqrt2$，出现于 $\omega = 1/RC$ 时：$f_c = 1/(2\pi RC) = \qty{1.0}{kHz}$。在
\qty{1.0}{kHz} 处：比值为 $0.71$，相位为 $-45^\circ$。在 \qty{10}{kHz} 处：
$RC\omega = 10$，比值为 $1/\sqrt{101} = 0.10$，相位为 $-84^\circ$。
\end{solution}

\begin{solution}{exo:b1:sinusoidal-impedance:6}
$\underline Z = 100 + j(250 - 100) = 100 + 150j$：$Z = \qty{180}{\ohm}$，
$\arg = 56^\circ$。$I = 1.0/180 = \qty{5.6}{mA}$，滞后于电源 $56^\circ$。
$U_R = \qty{0.56}{V}$，$U_L = \qty{1.39}{V}$，$U_C = \qty{0.56}{V}$；
相量图中：$U_R$ 沿 $I$ 方向，$U_L$ 向上，$U_C$ 向下，合成为 \qty{1.0}{V}。
$L\omega > 1/C\omega$：处于共振频率以上（感性）。
\end{solution}

\begin{solution}{exo:b1:sinusoidal-impedance:7}
$f_0 = 1/(2\pi\sqrt{LC})$：\qty{50}{pF} 时为 \qty{1.59}{MHz}，\qty{500}{pF} 时为
\qty{0.50}{MHz}。在 \qty{1.00}{MHz} 处：$C = 1/(4\pi^2f_0^2L) = \qty{127}{pF}$。
$\Delta f = \qty{10}{kHz}$。电台位于 \qty{1.009}{MHz}：$x = 1.009$，
$Q(x - 1/x) = 1.8$，比值为 $1/\sqrt{1 + 3.2} = 0.49$。电容器电压为
$QE = \qty{1.0}{mV}$。
\end{solution}

\begin{solution}{exo:b1:sinusoidal-impedance:8}
$Q(x - 1/x) = \pm1 \Leftrightarrow x^2 \mp x/Q - 1 = 0$；正根为
$x_{1,2} = \mp 1/2Q + \sqrt{1 + 1/4Q^2}$。二者之差为 $1/Q$，故
$\Delta\omega = \omega_0/Q$；乘积为 $(1 + 1/4Q^2) - 1/4Q^2 = 1$，故
$\omega_1\omega_2 = \omega_0^2$（通带在对数刻度上对称）。
\end{solution}

\begin{solution}{exo:b1:sinusoidal-impedance:9}
$I = P/(U\cos\varphi) = 2000/(230 \times 0.70) = \qty{12.4}{A}$。无功
功率为 $P\tan\varphi = 2000 \times 1.02 = \qty{2.0}{kvar}$。$C = 2040/(314
\times 230^2) = \qty{123}{\micro F}$。补偿后电流为 $2000/230 = \qty{8.7}{A}$。
损耗：补偿前为 $0.50 \times 12.4^2 = \qty{77}{W}$，补偿后为 $0.50 \times 8.7^2 =
\qty{38}{W}$。
\end{solution}

\begin{solution}{exo:b1:sinusoidal-impedance:10}
$\underline Y = 1/R + j(C\omega - 1/L\omega)$；$\underline U = \underline I/
\underline Y$ 在电纳为零时最大，即在 $\omega_0$ 处，此时 $U = RI$。半功率点满足
$\abs{C\omega - 1/L\omega} = 1/R$，即
$RC\omega_0(x - 1/x) = \pm1$：$\Delta\omega = \omega_0/(RC\omega_0) = 1/RC$，
$Q = RC\omega_0 = R\sqrt{C/L}$。数值：$\omega_0 = \qty{1.0e6}{rad/s}$
（$f_0 = \qty{159}{kHz}$），$Q = 10^4\sqrt{10^{-6}} = 10$，$\Delta f =
\qty{16}{kHz}$。
\end{solution}

\begin{solution}{exo:b1:sinusoidal-impedance:11}
$U_C = I/C\omega = E/[C\omega\sqrt{R^2 + (L\omega - 1/C\omega)^2}]$；
将根号内乘以 $C\omega$：$\sqrt{(RC\omega)^2 + (LC\omega^2 - 1)^2} =
\sqrt{x^2/Q^2 + (1 - x^2)^2}$。令 $s = x^2$，则 $(1 - s)^2 + s/Q^2$ 在
$s = 1 - 1/2Q^2$（若该值为正）时取最小值 $1/Q^2 - 1/4Q^4$，故
$U_{C,\max} = QE/\sqrt{1 - 1/4Q^2}$。$Q = 3.16$ 时：$x_r = 0.975$，
$U_{C,\max} = 3.2E$。当 $Q < 1/\sqrt2$ 时，最小值位于 $s = 0$，$U_C$
从 $E$ 起单调下降，无峰值。
\end{solution}

\begin{solution}{exo:b1:sinusoidal-impedance:12}
共振时 $i = I\cos\omega_0 t$，$u_C = (I/C\omega_0)\sin\omega_0 t$，
利用 $1/C\omega_0^2 = L$，
\[
\tfrac12 Li^2 + \tfrac12 Cu_C^2 = \tfrac12 LI^2\cos^2\omega_0 t
+ \tfrac12 LI^2\sin^2\omega_0 t = \tfrac12 LI^2 .
\]
每个周期耗散：$\tfrac12 RI^2T$。比值 $\times 2\pi$：$2\pi L/(RT)
= L\omega_0/R = Q$。
\end{solution}

\begin{solution}{pb:b1:sinusoidal-impedance:1}
\textbf{1.} $\underline Z = R + j(L\omega - 1/C\omega)$。

\textbf{2.} $\underline I = E/\underline Z$；$I = E/\sqrt{R^2 + (L\omega -
1/C\omega)^2}$；$\varphi_i = -\arctan[(L\omega - 1/C\omega)/R]$。

\textbf{3.} 电抗在 $\omega_0 = 1/\sqrt{LC}$ 处为零：$I_{\max}
= E/R$，$\varphi_i = 0$。

\textbf{4.} $C = 1/(4\pi^2f_0^2L) = \qty{127}{pF}$；$L\omega_0 = 2\times10^{-4}
\times 6.28\times10^6 = \qty{1257}{\ohm}$。

\textbf{5.} $Q = 1257/12.6 = 100$。

\textbf{6.} $L\omega - 1/C\omega = L\omega_0(x - 1/x)$，除以 $R$。

\textbf{7.} 当 $x < 1$ 时电抗为负，$\varphi_i > 0$，电流
超前（容性）；$x > 1$ 时则滞后（感性）。

\textbf{8.} $Q(x - 1/x) = \pm1$。

\textbf{9.} 由 $x^2 \mp x/Q - 1 = 0$，得 $x_{1,2} = \mp1/2Q + \sqrt{1 + 1/4Q^2}$，
$x_2 - x_1 = 1/Q$，故 $\Delta\omega = \omega_0/Q$。

\textbf{10.} $\Delta f = \qty{10}{kHz}$。

\textbf{11.} \qty{1.010}{MHz}处：$Q(x - 1/x) = 100 \times 0.0199 = 2.0$，
比值 $1/\sqrt{5.0} = 0.45$，功率 $0.20$：\qty{-7}{dB}。\qty{1.020}{MHz}处：
$3.96$，比值 $0.25$，功率 $0.06$：\qty{-12}{dB}。

\textbf{12.} 处于临界：邻近频道仅衰减 \qty{7}{dB}。提高
$Q$（即减小 $R$），或级联多个调谐级，可使
选择性更强。

\textbf{13.} \qty{0.50}{MHz} 至 \qty{1.59}{MHz}：调幅广播波段。

\textbf{14.} $\underline U_C = \underline I/jC\omega$；在 $\omega_0$ 处，
$U_C = I_{\max}/C\omega_0 = E/(RC\omega_0) = QE$。

\textbf{15.} $I_{\max} = 10^{-5}/12.6 = \qty{0.79}{\micro A}$；$U_C = QE =
\qty{1.0}{mV}$，无需放大器即可获得百倍的电压增益。

\textbf{16.} $\underline U_L = jL\omega_0\underline I = jQE$，$\underline U_C
= -jQE$：大小相等、方向相反，二者相互抵消；电源所见仅为 $R$。

\textbf{17.} 储存能量 $\tfrac12 LI_{\max}^2 = \qty{6.2e-17}{J}$；每周期耗散
能量 $\tfrac12 RI_{\max}^2T = \qty{3.9e-18}{J}$；$2\pi \times 6.2/0.39
\approx 100 = Q$。

\textbf{18.} $\tau = 2Q/\omega_0 = \qty{32}{\micro s}$。电路无法
跟随快于 $\sim\tau$ 的变化；带宽为 $\Delta f$ 的电路
只允许慢于 $\Delta f$ 的调制通过；取 $Q = 1000$，$\Delta f =
\qty{1}{kHz}$，节目中的高音将会丢失。

\textbf{19.} $Q$ 减半为 $50$，$\Delta f$ 加倍为 \qty{20}{kHz}，$U_C$
减半为 \qty{0.5}{mV}。

\textbf{20.} $R_{\text{tot}} = \qty{62.6}{\ohm}$：$Q = 1257/62.6 = 20$，
$\Delta f = \qty{50}{kHz}$，$U_C = 20 \times \qty{10}{\micro V} = \qty{0.2}{mV}$：
灵敏度与选择性均大幅下降。

\textbf{21.} $P = E^2/(2R_{\text{tot}}) = 10^{-10}/125 = \qty{8e-13}{W}$；
在 $R_a$ 上的占比：$50/62.6 = 80\%$。

\textbf{22.} 匹配要求负载电阻等于 $R_a$，即
总电阻要大；选择性则要求总电阻尽可能小
（$Q = L\omega_0/R_{\text{tot}}$）。采用直接连接时，
二者不可兼得。

\textbf{23.} 采用抽头或小耦合电容，使天线只接到谐振回路的
一部分：天线电阻经电路“看到”后大大减小，
从而保持 $Q$ 较高，同时天线
仍能在接近最佳值处输送功率。

\textbf{24.} 在半功率频率处，电抗等于 $\pm R$，
故 $\varphi = \pm45^\circ$，$\cos\varphi = 1/\sqrt2$；也有 $I =
I_{\max}/\sqrt2$，故 $P = \tfrac12 EI\cos\varphi$ 为其
谐振值的一半。

\textbf{25.} $Q$ 使信号放大 $Q$ 倍，并将带宽收窄至
$f_0/Q$，代价是响应变慢（$\tau = 2Q/\omega_0$）、与天线的
耦合变得脆弱；对于 \qty{10}{kHz} 的信道，
位于 \qty{1}{MHz} 处时，$Q \approx 100$ 即为所需数值——而天线必须
保持轻度耦合以维持该值。
\end{solution}
